\documentclass[10pt,a4paper]{amsart}
\usepackage[margin=19mm]{geometry}
\usepackage{amsmath,amssymb,mathtools}
\usepackage[scr=rsfs,cal=euler]{mathalfa}
\usepackage{xfrac}
\usepackage[unicode,hidelinks,bookmarks=false]{hyperref}
\newcommand{\ie}{i.e.\ }
\newcommand{\cf}{cf.\ }
\newcommand{\resp}{respectively\ }
\newcommand{\Subsection}{\subsection}
\providecommand{\nobreakdash}{\nobreak\hbox{-}}
\setlength{\emergencystretch}{2em}
\multlinegap0pt
\usepackage{amssymb}
\usepackage{mathtools}
\usepackage{enumerate}
\usepackage{xcolor}
\newtheorem{proposition}{Proposition}
\newcommand{\Q}{\mathbb{Q}}
\newcommand{\op}[1]{\operatorname{#1}}
\title{Hilbert's tenth problem for families of $ \mathbb{Z}_p $-extensions of imaginary quadratic fields}
\author{Katharina M\"uller, Anwesh Ray}
\date{Source: arXiv:2406.01443v2 (CC BY 4.0)}
\begin{document}
\maketitle

\noindent\emph{Source and licence.} Hilbert's tenth problem for families of $ \mathbb{Z}_p $-extensions of imaginary quadratic fields. Version arXiv:2406.01443v2 (CC BY 4.0), distributed by arXiv under the Creative Commons Attribution 4.0 International licence, http://creativecommons.org/licenses/by/4.0/, as recorded in the arXiv licence metadata for this exact version and shown on the abstract page https://arxiv.org/abs/2406.01443v2. Full licence notice: \texttt{licenses/arxiv-2406.01443-CC-BY-4.0.txt}.\par\smallskip

\noindent\textit{Editorial scope.} Counted unit: the article's proposition prop:selmer-decomposition-ord (source0.tex lines 303--308). The article's complete original proof of it is delivered with it (source0.tex lines 309--317, one \texttt{proof} environment of 9 lines). No mathematics is written, repaired or re-derived: the statement and the proof are the article's own lines, in the article's own notation, and no retained line is substituted or re-flowed.  No further source-local context is needed: every cross-reference of every retained line resolves inside the two delivered spans.  Nothing was omitted from any retained span, so the package carries no omission record.  No retained line contains a citation, so the wrapper carries no bibliography and no bibliography entry.  No retained line contains a figure environment, an image-inclusion command or drawing code, so no image is delivered and the package declares no asset dependency.  Editorial formatting only, in addition: an AMS \texttt{amsart} preamble that restates the article's own declarations and macros used by the retained lines, namely the environments \texttt{proposition} on separate counters, together with the standard packages those lines need.  The retained environments print in this document as Proposition 1 in order of appearance, whereas the article prints the same items with the numbering of its own full document.  The complete native source of the article as distributed in the arXiv e-print for this exact version is bundled unchanged as \texttt{sources/160206/source0.tex}.
\begin{proposition}\label{prop:selmer-decomposition-ord}
    Suppose $ E/\mathbb{Q} $ is an elliptic curve, and let $ E^{(K)} $ be the twist of $ E $ corresponding to $ K $. Then,
    \[
    \mathrm{Sel}_{p^\infty}(E/K_{\mathrm{cyc}}) \simeq \mathrm{Sel}_{p^\infty}(E/\mathbb{Q}_{\mathrm{cyc}}) \oplus \mathrm{Sel}_{p^\infty}(E^{(K)}/\mathbb{Q}_{\mathrm{cyc}}).
    \]
\end{proposition}

\begin{proof}
    From Shapiro's lemma we deduce that 
    \[H^1(\Q_\Sigma/K_{\op{cyc}}, E[p^\infty])\simeq H^1(\Q_\Sigma/\Q_{\op{cyc}}, E[p^\infty])\oplus H^1(\Q_\Sigma/\Q_{\op{cyc}}, E^{(K)}[p^\infty]),\] and 
    \[J_\ell(E/K_{\op{cyc}})=J_\ell(E/\Q_{\op{cyc}})\oplus J_\ell(E^{(K)}/\Q_{\op{cyc}}).\]
    These isomorphisms are compatible with the restriction maps, and we thus find that 
    \[
    \mathrm{Sel}_{p^\infty}(E/K_{\mathrm{cyc}}) \simeq \mathrm{Sel}_{p^\infty}(E/\mathbb{Q}_{\mathrm{cyc}}) \oplus \mathrm{Sel}_{p^\infty}(E^{(K)}/\mathbb{Q}_{\mathrm{cyc}}).
    \]
\end{proof}

\end{document}
